% ===========================================================================
%  comune/macro.tex — le abbreviazioni del progetto
% ===========================================================================
%
%  Servono a scrivere una cosa in un posto solo. Il nome di un pacchetto
%  Python compare duecento volte nella guida: se un giorno `approaches/`
%  diventa `overlays/`, si cambia qui e non duecento volte a mano — e
%  soprattutto non si finisce con metà documento che dice una cosa e metà
%  un'altra, che è il modo in cui la documentazione smette di essere creduta.
% ===========================================================================

% --- la lingua del documento -----------------------------------------------
%  \iflyraInglese la definisce comune/preambolo.tex. Qui la si crea se manca,
%  perché `make figure' include questo file senza il preambolo. Il nome del
%  condizionale si scrive con \csname: scritto per esteso dentro un ramo
%  saltato, TeX lo conterebbe come un \if aperto e sbaglierebbe i \fi.
\ifcsname iflyraInglese\endcsname\else
  \expandafter\newif\csname iflyraInglese\endcsname
\fi

% --- nomi del prodotto ------------------------------------------------------
\newcommand{\Lyra}{\textsc{Lyra}}
\newcommand{\LyraCapital}{\textsc{Lyra Capital}}
\newcommand{\app}{\iflyraInglese the application\else l'applicazione\fi}
\newcommand{\App}{\iflyraInglese The application\else L'applicazione\fi}
\newcommand{\sito}{\iflyraInglese the site\else il sito\fi}
\newcommand{\Sito}{\iflyraInglese The site\else Il sito\fi}

% --- i pacchetti del motore -------------------------------------------------
%  Si scrivono \pkgCore e non \file{core/}: così il giorno in cui una cartella
%  cambia nome, cambia qui.
\DeclareRobustCommand{\pkg}[1]{{\ttfamily\small\color{lyraInchiostro!78}#1}}
\newcommand{\pkgCore}{\pkg{core/}}
\newcommand{\pkgStrategie}{\pkg{strategies/}}
\newcommand{\pkgApprocci}{\pkg{approaches/}}
\newcommand{\pkgPesi}{\pkg{weights/}}
\newcommand{\pkgOverlay}{\pkg{overlays/}}
\newcommand{\pkgRegime}{\pkg{regime/}}
\newcommand{\pkgDati}{\pkg{data/}}
\newcommand{\pkgBacktest}{\pkg{backtest/}}
\newcommand{\pkgLive}{\pkg{live/}}
\newcommand{\pkgEsecuzione}{\pkg{execution/}}
\newcommand{\pkgReport}{\pkg{reporting/}}
\newcommand{\pkgDashboard}{\pkg{dashboard/}}
\newcommand{\pkgDominio}{\pkg{domain/}}
\newcommand{\pkgPorte}{\pkg{ports/}}
\newcommand{\pkgApplicazione}{\pkg{application/}}
\newcommand{\pkgAdattatori}{\pkg{adapters/}}
\newcommand{\pkgComposizione}{\pkg{composition/}}
\newcommand{\pkgApi}{\pkg{api/}}

% --- termini del dominio, con la prima occorrenza indicizzata ---------------
\newcommand{\term}[1]{\emph{#1}\index{#1}}
\newcommand{\termdef}[1]{\textbf{#1}\index{#1}}

% --- riferimenti a classi, funzioni, file -----------------------------------
%  \simbolo{Nome}            una classe o una funzione
%  \sorgente{file.py}{120}   un punto preciso del sorgente
\DeclareRobustCommand{\simbolo}[1]{\cod{#1}}
\newcommand{\sorgente}[2]{{\ttfamily\footnotesize\color{lyraInchiostro!60}#1:#2}}

% --- tasti e voci dell'interfaccia, per il manuale utente -------------------
%  \voce{Overview}           una voce di menu o una scheda
%  \pulsante{Run backtest}   un pulsante da premere
%  \campo{Target vol}        un campo da compilare
%  \percorso{Settings > Data} un percorso di navigazione
\newcommand{\voce}[1]{\textsf{\textbf{#1}}}
\newcommand{\pulsante}[1]{%
  \tcbox[on line, boxsep = 1pt, left = 3pt, right = 3pt, top = 0.5pt,
         bottom = 0.5pt, arc = 3pt, boxrule = 0.5pt,
         colback = lyraInchiostro, colframe = lyraInchiostro,
         coltext = white]{\sffamily\scriptsize\bfseries #1}}
\newcommand{\pulsantesec}[1]{%
  \tcbox[on line, boxsep = 1pt, left = 3pt, right = 3pt, top = 0.5pt,
         bottom = 0.5pt, arc = 3pt, boxrule = 0.5pt,
         colback = white, colframe = lyraInchiostro!50,
         coltext = lyraInchiostro]{\sffamily\scriptsize\bfseries #1}}
\newcommand{\campo}[1]{\textsf{\itshape #1}}
\newcommand{\percorso}[1]{\textsf{\small #1}}
\newcommand{\freccia}{\;\textcolor{lyraAccento}{$\rightarrow$}\;}

% --- una rotta HTTP ---------------------------------------------------------
%  \rotta{GET}{/api/marketplace/<id>/perche-non-mio}
%
%  Esiste perche' \cod{} e' una scatola indivisibile: in una colonna di
%  tabella larga cinque centimetri un percorso lungo sborda, e non c'e'
%  niente che TeX possa fare. Qui il percorso e' testo normale e puo' andare
%  a capo dopo ogni barra, che e' il punto naturale.
%
%  Il ciclo prende un carattere alla volta: funziona perche' un percorso non
%  contiene spazi (un argomento non delimitato li mangerebbe) ed e' in ASCII
%  (una lettera accentata e' un carattere attivo e manda in errore il ciclo).
%  Per qualunque altra cosa si usa \cod{}.
\makeatletter
\def\lyra@finePercorso{\lyra@finePercorsoMARK}
\def\lyra@percorso#1{%
  \ifx#1\lyra@finePercorso
    \let\lyra@prossimo\relax
  \else
    #1%
    % Dopo questi quattro caratteri la riga puo' andare a capo. Sono i punti
    % in cui un percorso o un nome di file si leggono comunque bene spezzati.
    \ifx#1/\allowbreak\else
      \ifx#1\_\allowbreak\else
        \ifx#1.\allowbreak\else
          \ifx#1-\allowbreak\fi\fi\fi\fi
    \let\lyra@prossimo\lyra@percorso
  \fi
  \lyra@prossimo}
\DeclareRobustCommand{\percorsoHTTP}[1]{%
  {\ttfamily\small\color{lyraInchiostro!85}%
   \lyra@percorso#1\lyra@finePercorso}}
\makeatother

\DeclareRobustCommand{\rotta}[2]{%
  {\sffamily\scriptsize\bfseries\color{lyraAccentoScuro}#1}\,\percorsoHTTP{#2}}

% --- la figura standard -----------------------------------------------------
%  Centra, mette la didascalia sotto e l'etichetta. Esiste perché le tre righe
%  scritte a mano ogni volta divergono: una figura senza \label non si può
%  citare, e una citata per numero si rompe appena se ne aggiunge un'altra.
%  #1 file da includere (senza estensione)  #2 didascalia  #3 etichetta
%  `max width` e non `width`: una figura piccola resta della sua misura, una
%  piu' larga della giustezza viene rimpicciolita invece di sbordare nel
%  margine. Senza, ogni figura andrebbe dimensionata a mano e ridimensionata
%  a ogni modifica del testo.
\newcommand{\figuraTikz}[3]{%
  \begin{figure}[htbp]
    \centering
    \begin{adjustbox}{max width = \linewidth}
      \input{figure/tikz/#1}
    \end{adjustbox}
    \caption{#2}\label{fig:#3}
  \end{figure}}

%  Come sopra ma per un'immagine, con larghezza.
%  #1 larghezza  #2 file  #3 didascalia  #4 etichetta
\newcommand{\figuraImg}[4]{%
  \begin{figure}[htbp]
    \centering
    \includegraphics[width = #1\linewidth]{#2}
    \caption{#3}\label{fig:#4}
  \end{figure}}

%  Schermata con cornice: una schermata senza bordo si confonde con il fondo
%  della pagina, e non si capisce dove finisce l'interfaccia. La cornice sta
%  DENTRO la larghezza chiesta: altrimenti una schermata larga quanto la
%  pagina sborda di un punto per lato.
\newcommand{\figuraSchermata}[4]{%
  \begin{figure}[htbp]
    \centering
    \setlength{\fboxsep}{0pt}\setlength{\fboxrule}{0.5pt}%
    {\color{lyraGrigioChiaro}\fbox{\includegraphics[width = \dimexpr #1\linewidth - 2\fboxrule\relax]{#2}}}
    \caption{#3}\label{fig:#4}
  \end{figure}}

% --- riferimenti incrociati con la parola davanti ---------------------------
%  In inglese la parola vuole la maiuscola («see Chapter 4»), in italiano no.
\newcommand{\vediFig}[1]{\iflyraInglese Figure\else figura\fi~\ref{fig:#1}}
\newcommand{\vediTab}[1]{\iflyraInglese Table\else tabella\fi~\ref{tab:#1}}
\newcommand{\vediCap}[1]{\iflyraInglese Chapter\else capitolo\fi~\ref{cap:#1}}
\newcommand{\vediSez}[1]{\iflyraInglese Section~\else \S\,\fi\ref{sez:#1}}

% --- numeri e unità ---------------------------------------------------------
%  Ogni lingua con la sua convenzione, la stessa dell'applicazione e del sito:
%  in italiano «10\,%» e «49\,€», in inglese «10%» e «€49» (le migliaia, in
%  inglese, si scrivono con la virgola nell'argomento: \eur{1,490}).
\newcommand{\perc}[1]{\iflyraInglese #1\%\else #1\,\%\fi}
\newcommand{\eur}[1]{\iflyraInglese \texteuro#1\else #1\,\texteuro\fi}
